\subsection{下降态射}

设 X 和 Y 是拓扑空间，f 是从 X 到 Y 的连续映射。设 \((Z,p)\) 和 \((Z',p')\) 是赋予关于 f 的下降数据的 X-空间，分别记为 \(\tau\) 和 \(\tau'\)。称 X-态射 \(\varphi\colon Z\to Z'\) 关于下降数据 \(\tau\) 和 \(\tau'\) 相容，是指有

\[
\tau^ {\prime} (\varphi (z), x) = \varphi (\tau (z, x))
\]

对所有 \((z,x)\in\mathbf{Z}\times_{\mathbf{Y}}\mathbf{X}\)。等价地说，\(\varphi\) 的像将关于关系 \(R_{\tau}\) 等价的两个点变为关于关系 \(R_{\tau^{\prime}}\) 等价的点。这样的态射 \(\varphi\) 通过取商定义出连续映射 \(\overline{\varphi}:Z/R_{\tau}\to Z^{\prime}/R_{\tau^{\prime}}\)；它是 Y-空间的态射。

记 \(\mathcal{C}_{\tau,\tau^{\prime}}(\mathrm{Z};\mathrm{Z}^{\prime})\) 为从 Z 到 \(Z^{\prime}\) 且与下降数据 \(\tau\) 和 \(\tau^{\prime}\) 相容的 X-态射集合。

命题 3. --- 设 X 和 Y 是拓扑空间，设 \(f\colon X \to Y\) 是连续映射。设 \((Z,p)\) 和 \((Z',p')\) 是赋予关于 f 的下降数据的 X-空间，分别记为 \(\tau\) 和 \(\tau'\)。若下降数据 \(\tau'\) 有效，则映射 \(\varphi \mapsto \overline{\varphi}\) 是从 \(\mathcal{C}_{\tau,\tau'}(Z;Z')\) 到 \(\mathcal{C}_Y(Z/R_\tau;Z'/R_{\tau'})\) 的双射。

对于每个 X-态射 \(\varphi\)（从 Z 到 \(Z'\) 且与下降数据相容），映射 \(\overline{\varphi}\) 是 Y-态射。反之，设 \(g\colon Z\to Z/R_{\tau}\) 和 \(g'\colon Z'\to Z'/R_{\tau'}\) 是规范映射，设 \(\psi\colon Z/R_{\tau}\to Z'/R_{\tau'}\) 是 Y-态射。映射 \(p\colon Z\to X\) 和 \(\psi\circ g\colon Z\to Z'/R_{\tau'}\) 都是 Y-态射。下降数据 \(\tau'\) 有效这一假设意味着图表

\[
\begin{array}{c} \mathrm {Z^ {\prime}} \xrightarrow {g ^ {\prime}} \mathrm {Z}^ {\prime} / \mathrm {R}_ {\tau^ {\prime}} \\ p ^ {\prime} \Biggl \downarrow \qquad \qquad \qquad \qquad \Biggl \downarrow q ^ {\prime} \\ \mathrm{X} \xrightarrow {f} \mathrm{Y} \end{array}
\]

是笛卡尔方形。因此存在唯一连续映射 \(\varphi\colon Z\to Z'\)，使得 \(p'\circ\varphi=p\) 且 \(g'\circ\varphi=\psi\circ g\)。第一个等式意味着 \(\varphi\) 是 X-态射，第二个等式意味着 \(\varphi\) 与下降数据相容且 \(\overline{\varphi}=\psi\)，命题得证。

